\chapter{内核锁、死锁与可证明的等待}

\section{原理}

操作系统进入并发世界以后，错误不再只来自单个函数写错变量，而来自两个执行流同时触碰同一份状态。中断处理程序可能和普通内核线程同时访问设备队列；两个 CPU 可能同时把任务放入 runqueue；一个进程关闭 fd 时，另一个线程可能正在通过同一个 open file description 发起 I/O。锁的目的不是让代码“看起来串行”，而是给共享对象规定一个可证明的临界区：谁能读，谁能写，写入完成前谁必须等待，等待时是否允许睡眠。

从裸机单片机演进过来时，最早的“锁”通常是关中断。它简单有效，因为只有一个 CPU，且并发主要来自中断打断主循环。进入多核后，关本地中断不能阻止另一个 CPU 访问同一对象，于是需要自旋锁；进入可能长时间等待的路径后，自旋浪费 CPU，于是需要 mutex、semaphore 和等待队列。每一种同步原语都绑定一个上下文限制：中断上下文不能睡眠；持有自旋锁时不能调用可能阻塞的函数；持有文件系统锁时不能随意回调用户内存复制路径；调度器自己的锁必须在切换前后保持严格顺序。

\begin{keyidea}
锁不是性能技巧，而是状态所有权的边界。设计锁时先写对象不变量，再决定哪些操作必须在同一把锁下完成。
\end{keyidea}

内核锁设计要同时回答四个问题：保护哪个对象，持锁期间允许做什么，等待者在哪里排队，失败路径如何释放。很多死锁来自第四个问题被忽略：正常路径释放锁，错误路径提前返回却忘了释放；或者关闭路径拿到 A 锁后等待 B，而另一个路径拿到 B 后等待 A。

\section{实践}

下面是常见内核锁原语的实践边界：

\begin{longtable}{p{0.18\linewidth}p{0.26\linewidth}p{0.24\linewidth}p{0.22\linewidth}}
\toprule
原语 & 适用场景 & 禁止事项 & 典型对象 \\
\midrule
关中断 & 单核短临界区，中断共享变量 & 多核共享保护，长时间计算 & tick 计数、单核设备寄存器影子状态 \\
自旋锁 & 多核短临界区，不能睡眠的上下文 & 持锁 I/O、持锁 page fault、持锁等待用户事件 & runqueue、描述符环、引用计数 \\
mutex & 进程上下文中可能等待的互斥 & 中断上下文获取，持有后进入不可重入路径 & inode 元数据、文件系统目录更新 \\
读写锁 & 读多写少且临界区短 & 长读者导致写者饥饿，升级锁顺序混乱 & mount 表、路由表快照 \\
RCU & 读路径极热，读者不阻塞写者 & 立即释放旧对象，读区内睡眠 & fdtable、路径缓存、配置表 \\
semaphore & 计数资源，允许多个持有者 & 用它表达复杂对象不变量 & buffer 数量、并发 I/O 限额 \\
\bottomrule
\end{longtable}

一个合理的锁设计文档至少包含下面字段：

\begin{lstlisting}
Object:
  runqueue[cpu]

Invariant:
  task.state == RUNNABLE implies task is on exactly one runqueue
  rq.nr_running == number of tasks linked from rq.head

Lock:
  rq.lock protects rq.head, rq.nr_running, task.rq_link

Allowed while holding:
  enqueue/dequeue, choose_next, update vruntime

Forbidden while holding:
  page allocation with reclaim, disk I/O, copy_from_user, waiting on mutex
\end{lstlisting}

等待队列必须把“检查条件”和“入队睡眠”放在同一把锁保护下，否则会出现 lost wakeup。正确模型不是“先看条件，不满足就睡”，而是“在保护条件的锁内，把自己放到等待集合中，然后原子地释放锁并切换状态”。唤醒方也在同一把锁下改变条件并唤醒等待者。

\section{算法与实现分析}

\subsection{自旋锁}

最小自旋锁可以用原子交换实现：

\begin{lstlisting}
lock(spinlock *l):
  while atomic_exchange(l->held, true) == true:
    cpu_relax()

unlock(spinlock *l):
  atomic_store_release(l->held, false)
\end{lstlisting}

这个算法的互斥性来自原子交换：同一时刻只有一个 CPU 能看到旧值为 false。复杂度不是固定常数，因为竞争时等待时间取决于持锁者运行时间和 CPU 调度。实现上必须使用 acquire/release 语义，否则 CPU 或编译器可能把临界区内存访问重排到锁外。真实内核还会加入持有者记录、递归检查、抢占关闭、IRQ 保存和等待统计。

普通 test-and-set 自旋锁在高竞争下会让所有 CPU 反复写同一个 cache line。ticket lock 用取号保证公平：

\begin{lstlisting}
lock(ticket_lock *l):
  my = fetch_add(l->next, 1)
  while load_acquire(l->owner) != my:
    cpu_relax()

unlock(ticket_lock *l):
  store_release(l->owner, l->owner + 1)
\end{lstlisting}

ticket lock 的优点是 FIFO 公平，缺点是所有等待者仍然轮询 owner cache line。队列锁进一步让每个等待者轮询自己的节点，降低缓存抖动。教学内核可以先实现 test-and-set，再用压力测试观察竞争成本，最后引入 ticket 或 MCS 解释公平性和 cache line 迁移。

\subsection{mutex 和等待队列}

mutex 的关键不是把自旋改成睡眠，而是在失败时把当前线程从 runnable 集合移到等待队列：

\begin{lstlisting}
mutex_lock(m):
  disable_preempt()
  spin_lock(m.wait_lock)
  if m.owner == null:
    m.owner = current
    spin_unlock(m.wait_lock)
    enable_preempt()
    return
  enqueue(m.waiters, current)
  current.state = BLOCKED
  schedule_unlocking(m.wait_lock)
  enable_preempt()
  retry mutex_lock(m)

mutex_unlock(m):
  spin_lock(m.wait_lock)
  if empty(m.waiters):
    m.owner = null
  else:
    next = dequeue(m.waiters)
    m.owner = next
    next.state = RUNNABLE
    enqueue_runqueue(next)
  spin_unlock(m.wait_lock)
\end{lstlisting}

这里有三个不变量：owner 要么为空要么指向持有线程；等待队列里的线程必须是 BLOCKED；被唤醒线程必须重新进入 runnable 集合。若 unlock 先释放 owner 再唤醒，另一个新线程可能插队，造成等待者饥饿。若唤醒时没有持有等待队列锁，等待节点可能和 timeout/cancel 路径并发释放。

\subsection{死锁条件与锁顺序}

经典死锁需要四个条件同时成立：互斥、持有并等待、不可抢占、循环等待。内核无法完全消除前三个条件，因为对象确实需要互斥，很多资源不能强制抢占。因此工程上主要打破循环等待：给锁分层，规定全局顺序。

\begin{lstlisting}
Allowed lock order:
  tasklist_lock
    -> mm.lock
      -> vma.lock
        -> page.lock
  superblock.lock
    -> inode.lock
      -> pagecache.lock
        -> buffer.lock
  netns.lock
    -> socket.lock
      -> skb_queue.lock
\end{lstlisting}

锁顺序不是文档装饰，而要能被测试和运行时检查。简化 lockdep 可以在每个线程记录当前持有锁栈，每次获取新锁时加入一条“已持有锁 -> 新锁”的有向边；若图出现环，就报告潜在死锁。

\begin{lstlisting}
on_lock_acquire(new_lock):
  for held in current.held_locks:
    graph_add_edge(held.class, new_lock.class)
    if graph_has_cycle():
      report_deadlock_risk(current, held, new_lock)
  push(current.held_locks, new_lock)
\end{lstlisting}

这个检查的成本取决于锁类数量和边数量，不能在极热路径无节制开启，但作为调试构建非常有价值。更重要的是，它把“某次测试没有卡死”升级成“锁顺序图没有出现已知环”。

\subsection{RCU 的读写分离}

RCU 解决的是读路径极热、写路径较少的场景。读者进入 RCU 临界区时不获取传统互斥锁，只保证自己不会在临界区内被回收对象悬空；写者发布新对象后，必须等所有旧读者经过 quiescent state 才释放旧对象。

\begin{lstlisting}
reader():
  rcu_read_lock()
  p = rcu_dereference(global_ptr)
  use(p)
  rcu_read_unlock()

writer():
  old = global_ptr
  new = clone_and_update(old)
  rcu_assign_pointer(global_ptr, new)
  call_rcu(old, free_after_grace_period)
\end{lstlisting}

RCU 的不变量是：任何读者看到的对象在读侧临界区结束前仍然有效。它不保证读者看到最新状态，也不适合读区内睡眠的普通实现。把 RCU 当成“无锁万能替代品”会导致非常隐蔽的 use-after-free。

\section{测试}

锁测试不能只跑正常路径。最低测试包括：

\begin{enumerate}
\tightlist
\item 互斥测试：多个 CPU 同时递增共享计数器，最终值必须等于操作次数。
\item lost wakeup 测试：让唤醒和入睡交错，等待者不能永久睡眠。
\item 中断上下文测试：持有自旋锁时触发同类中断，检查是否需要 irqsave。
\item 死锁检测测试：故意构造 A 后 B 与 B 后 A，lockdep 必须报告环。
\item 饥饿测试：大量短任务竞争 mutex，等待时间不能无限增长。
\item 退出路径测试：加锁后在每个错误分支注入失败，所有锁计数必须归零。
\end{enumerate}

\begin{rubricbox}
本章合格标准：能说明每把锁保护哪个对象；能判断自旋锁、mutex、RCU 的上下文限制；能写出 lost wakeup 的反例和修复；能用锁顺序图解释死锁风险。
\end{rubricbox}

\section{源码级补充：qspinlock、rwsem、seqlock 与 lockdep}

\subsection{raw spinlock、ticket lock 与 queued spinlock}

简单 test-and-set 锁在竞争下会让所有 CPU 反复写同一 cache line。ticket lock 提供 FIFO 公平，但所有等待者仍然轮询同一个 owner。queued spinlock 把等待者组织成队列，思想接近 MCS：每个 CPU 主要等待自己的节点，减少 cache line 抖动。

\begin{lstlisting}
mcs_lock(lock, node):
  node.locked = true
  pred = xchg(lock.tail, node)
  if pred != null:
    pred.next = node
    while node.locked:
      cpu_relax()

mcs_unlock(lock, node):
  if node.next == null:
    if cmpxchg(lock.tail, node, null):
      return
    wait until node.next != null
  node.next.locked = false
\end{lstlisting}

Linux 的 qspinlock 比这段伪代码更紧凑，读 \filepath{kernel/locking/qspinlock.c} 时重点看三件事：pending bit 如何处理短竞争，MCS queue 如何承载长竞争，锁路径如何配合 preempt 和 paravirt。不要把 queued lock 理解成“更复杂所以一定更快”，它是为高竞争多核场景付出的复杂度。

\subsection{rwlock、rwsem 和乐观自旋}

读写锁允许多个读者并行，一个写者独占。rwlock 通常用于短临界区，不能睡眠；rwsem 是可睡眠读写信号量，适合 VMA、文件系统等较长路径。rwsem 还可能做 optimistic spinning：如果持有者正在 CPU 上运行，等待者短暂自旋，期望持有者很快释放，避免睡眠/唤醒成本。

\begin{longtable}{p{0.18\linewidth}p{0.34\linewidth}p{0.34\linewidth}}
\toprule
原语 & 优点 & 风险 \\
\midrule
rwlock & 读者并行，路径短 & 长读者或写频繁时收益差，不能睡眠 \\
rwsem & 可睡眠，适合较长临界区 & 饥饿、公平性、唤醒策略复杂 \\
optimistic spin & 避免短等待睡眠成本 & 持有者不运行时浪费 CPU \\
\bottomrule
\end{longtable}

读源码可看 \filepath{kernel/locking/rwsem.c}。关注 wait list、owner、reader count、writer handoff。若一个路径持有 rwsem 后等待另一个也需要此 rwsem 的工作完成，就容易形成隐藏死锁。

\subsection{seqlock：写者优势的读一致性}

seqlock 用一个递增序列号表达写入是否发生。读者开始时读取 seq，复制数据，结束时再读 seq；如果 seq 变化或为奇数，重试。写者持锁并把 seq 变奇数，写完后变偶数。

\begin{lstlisting}
read:
  do {
    seq = read_seqbegin(&s)
    local = shared_data
  } while (read_seqretry(&s, seq))

write:
  write_seqlock(&s)
  shared_data = new_value
  write_sequnlock(&s)
\end{lstlisting}

seqlock 适合读多写少、读者可重试、数据可复制的场景，例如时间数据。它不适合读者不能重试、数据包含指针且对象生命周期复杂、写者频繁导致读者长期重试的场景。

\subsection{per-CPU 与 preempt/migrate 控制}

per-CPU 变量让每个 CPU 有自己的副本，减少锁和 cache line bouncing。但访问 per-CPU 数据时，要确保当前执行不会迁移到另一个 CPU，否则先取到 CPU0 的指针，随后迁移到 CPU1 继续使用，会破坏语义。

\begin{lstlisting}
this_cpu_counter_add(delta):
  preempt_disable()
  per_cpu(counter, smp_processor_id()) += delta
  preempt_enable()
\end{lstlisting}

有些路径用 \code{migrate\_disable} 而不是完全禁止抢占，表示任务可以被抢占但不能迁移 CPU。读调度和网络栈源码时，这类边界经常出现。per-CPU 不是无锁魔法，它把共享变成分片，同时引入汇总、CPU hotplug 和迁移限制。

\subsection{内存屏障原语}

\begin{longtable}{p{0.22\linewidth}p{0.52\linewidth}}
\toprule
原语 & 语义 \\
\midrule
\code{smp\_rmb} & 约束读-读顺序 \\
\code{smp\_wmb} & 约束写-写顺序 \\
\code{smp\_mb} & 完整内存屏障，约束读写重排 \\
\code{smp\_store\_release} & 发布数据前的写入对获取方可见 \\
\code{smp\_load\_acquire} & 看到发布指针后能看到发布前初始化 \\
\bottomrule
\end{longtable}

屏障必须服务于具体协议。若无法说出“谁发布、谁获取、保护哪组字段”，就不应随便加 \code{smp\_mb}。错误屏障可能既没修 bug，又增加维护负担。

\subsection{lockdep 报告怎么读}

lockdep 报告通常会给出两个锁类、已观察到的获取顺序、当前试图形成的新顺序和调用栈。读报告时先找锁类名，再画边：A 后 B，B 后 A 是否形成环；然后检查是否真实同一实例、同一上下文，还是锁类标注过粗。

\begin{lstlisting}
Observed order:
  inode_lock -> page_lock
New order:
  page_lock -> inode_lock
Result:
  possible circular locking dependency
\end{lstlisting}

不要把 lockdep 当成“测试失败噪声”。它发现的是潜在等待环，可能在当前负载下未死锁，但在错误时序下会卡死。高质量内核代码要能解释每条锁边的必要性，或调整锁类/顺序消除环。
